\chapter{Simulatore hardware C++}

Il file \texttt{hardware/main.cpp} implementa la classe \texttt{PlantPulseDevice} e il ciclo principale del simulatore.

\section{Misurazioni}
Il dispositivo simula quattro parametri: umidità del terreno, luminosità, temperatura dell'aria e umidità dell'aria. I valori iniziali sono impostati nel codice e vengono modificati gradualmente con piccole variazioni casuali; i valori sono limitati a intervalli plausibili definiti nel simulatore.

Il metodo che prepara la misurazione costruisce un JSON con \texttt{device\_id}, timestamp UTC e valori numerici. La funzione di invio utilizza libcurl per effettuare una POST verso \texttt{http://127.0.0.1:8000/measurements}.

\section{Gestione dei comandi}
Il simulatore recupera i comandi pendenti tramite l'endpoint dedicato. Sono gestiti i comandi per impostare l'intervallo di misurazione e i quattro parametri ambientali. Dopo l'applicazione, il simulatore invia un acknowledgement con l'esito.

\section{Ciclo operativo}
Il ciclo principale aggiorna i valori simulati, controlla i comandi e invia le misurazioni secondo l'intervallo configurato. L'intervallo iniziale è di due secondi. Il programma resta attivo fino all'interruzione del processo.
